行业周期仍强，但已经从“谁能抢到更多订单”转向“谁能把高价订单按时、按成本交出来”。2025年中国新接订单同比回落，而手持和完工继续增长，说明流量高点可能已过、存量兑现仍在加速。对600150而言，行业景气只能解释外部环境；真正决定估值的是2023--2024年订单在2026--2028年的价格、结构、成本和现金兑现。

\section{E06｜全球周期：高在手与高船价并存，新签边际回落}

全球2025年新签约约5,800万CGT、1,860亿美元，CGT虽同比下降，仍较十年均值高约30\%；交付约4,380万CGT，同比增加约6\%。OECD口径全球在手约1.69亿CGT，对当年交付的机械比约3.86倍，只能说明行业工作量较长，不能替代公司交期。\srcid{II-08, II-10}

\begin{exhibitbox}[全球造船周期仪表盘]
\small
\begin{tabularx}{\exhibitboxwidth}{L{3.0cm}R{3.0cm}L{4.2cm}X}
\toprule
\textbf{指标} & \textbf{最新观察} & \textbf{周期含义} & \textbf{600150含义} \\
\midrule
全球新签合同 & 约5,800万CGT/1,860亿美元 & 较2024高位回落，价值量仍高 & 新单不宜按高峰线性外推 \\
全球交付 & 约4,380万CGT，+6\% & 行业由签单转向生产兑现 & 设备、人工与项目管理成为瓶颈 \\
全球在手 & 约1.69亿CGT & 船位和工作量仍偏紧 & 支撑议价，但增加延期与组织风险 \\
新船价格指数 & 约184 & 仍处历史高位 & 旧订单批次逐步体现价格优势 \\
公开交期观察 & 约3.7年 & 优质船位稀缺 & 交付纪律比订单数量更重要 \\
\bottomrule
\end{tabularx}
\sourcenote{II-08、II-10及中船协公开行业材料；全球数据不证明600150公司份额。}
\end{exhibitbox}

周期的“所以呢”是：高在手降低近期收入断崖概率，却提高市场对准时交付的要求。若交付延迟、设备涨价或返工侵蚀毛利，高船价并不会自动转化为高利润；若新签持续回落，估值还会提前反映2028年后的正常化。

\section{E07｜中国份额：供应优势明确，但不是600150份额}

工信部披露，2025年中国造船完工5,369万DWT、同比+11.4\%，新接10,782万DWT、同比-4.6\%，手持27,442万DWT、同比+31.5\%；相应全球份额为56.1\%、69.0\%和66.8\%。这些是国家DWT口径，不能推导600150价值份额或公司CR3/CR5。\srcid{II-06}

\begin{exhibitbox}[中国造船三大指标]
\centering
\begin{tikzpicture}
\begin{axis}[
  ybar,
  width=0.93\exhibitwidth,
  height=6.4cm,
  bar width=13pt,
  ymin=0,ymax=30000,
  ylabel={万DWT},
  symbolic x coords={完工,新接,手持},
  xtick=data,
  axis line style={draw=rulegrey},
  tick style={draw=rulegrey},
  ymajorgrids=true,
  grid style={draw=panelgrey},
  nodes near coords,
  nodes near coords style={font=\scriptsize},
  every axis plot/.append style={fill=accentblue!75,draw=accentblue}
]
\addplot coordinates {(完工,5369) (新接,10782) (手持,27442)};
\end{axis}
\end{tikzpicture}
\sourcenote{工信部II-06。国家DWT数据仅作行业背景。}
\end{exhibitbox}

国家供应优势有两面：一方面，中国船位、供应链和工程人才形成规模效应；另一方面，新增产能与民营船厂效率也会在订单流量回落时加剧竞争。600150的2025新签DWT相当于中国总量约28.3\%，只是体量代理，船型、海工和价值结构不同，不能称为价值市占率。

\section{需求锚点只解释“为何下单”}

集装箱班轮、能源运输、散货更新、汽车出口、邮轮与海工资本开支形成不同需求锚。IMO减排框架支持长期绿色更新，但截至截止日尚未形成生效规则；政策方向不能直接资本化为双燃料船收入。国防需求具战略属性，但订单和经济性未披露，不进入军品SOTP。\srcid{II-12, CF-03}

\begin{exhibitbox}[需求锚、船型与主要脆弱点]
\small
\begin{tabularx}{\exhibitboxwidth}{L{2.8cm}L{3.4cm}L{3.4cm}X}
\toprule
\textbf{需求锚} & \textbf{主要船型} & \textbf{正向逻辑} & \textbf{脆弱点} \\
\midrule
班轮资本开支 & 大型/双燃料集装箱船 & 单船价值和复杂度较高 & 班轮盈利与资本开支周期 \\
能源贸易重构 & 油轮、LNG/LPG/VLEC等 & 船龄更新、航线重构与技术门槛 & 客户融资、围护/燃料系统许可 \\
干散货更新 & Capesize、VLOC、Kamsarmax & 老龄替换与标准化生产 & 标准船竞争强、价格透明 \\
汽车出口 & PCTC/RoRo & 复杂度与单船价值较高 & 贸易政策与车运供需 \\
邮轮与文旅 & 大型邮轮/客滚 & 复杂系统集成与系列化学习 & 国产化约35\%，外部系统依赖 \\
海洋能源 & FPSO/平台/支持船 & 高价值项目制机会 & 减值、延期和验收风险高 \\
\bottomrule
\end{tabularx}
\sourcenote{II-07--II-12、CF-03；需求锚不证明600150已获得特定订单。}
\end{exhibitbox}

对估值而言，需求锚只决定未来订单池的可能方向；公司收入仍需“订单归属\(\rightarrow\)船厂排程\(\rightarrow\)履约确认”的证据链。绿色、高端或军工标签若缺少公司合同额、交期和毛利，不进入基础或牛市价值。

\section{周期位置与估值持续期}

2026年更像利润率验证年，而不是订单方向验证年。公司计划交付168艘、1,650万DWT，主要承接2023--2024年订单；若高船价与结构升级开始兑现，2026利润会显著上升。但市场定价更依赖2027--2028：新签流量是否足以补充高价值订单，设备与人工成本是否抬升，盈利能否转为现金。

\begin{keyinsight}[周期判断的行动含义]
强周期并不自动等于股票低估。当前价格已经承认2026利润上行，新增回报需要证明高利润持续期比市场预期更长。应重点跟踪新签价值结构、交付兑现、核心毛利和现金，而不是重复引用国家份额与全球TAM。
\end{keyinsight}

周期风险将在两条路径上进入估值：已签订单通过延期、成本和现金影响2026--2027盈利；未来订单流通过价格、覆盖年限和终值影响2028年后倍数。两者不能混为“造船周期向上”一句话。
